% FCS 详细技术文档 — XeLaTeX
% 编译：cd docs/latex && make
\documentclass[UTF8,a4paper,11pt,fancyhdr,oneside]{ctexart}

\usepackage{geometry}
\geometry{margin=2.2cm,headheight=14pt}
\usepackage{fontspec}
\setCJKmainfont{Noto Serif CJK SC}
\setCJKsansfont{Noto Sans CJK SC}
\setCJKmonofont{Noto Sans Mono CJK SC}
\setmainfont{DejaVu Serif}
\setsansfont{DejaVu Sans}
\setmonofont{DejaVu Sans Mono}

\usepackage{booktabs}
\usepackage{longtable}
\usepackage{tabularx}
\usepackage{multirow}
\usepackage{array}
\usepackage{graphicx}
\usepackage{xcolor}
\usepackage{colortbl}
\usepackage{hyperref}
\usepackage{enumitem}
\usepackage{amsmath,amssymb}
\usepackage{listings}
\usepackage{caption}
\usepackage{subcaption}
\usepackage{tikz}
\usetikzlibrary{arrows.meta,positioning,fit,calc,shapes.geometric,shapes.multipart,chains,matrix,decorations.pathreplacing}
\usepackage{fancyhdr}
\usepackage{microtype}
\usepackage{algorithm}
\usepackage{algpseudocode}
\usepackage{float}
\usepackage{mdframed}

\definecolor{oldbg}{RGB}{255,240,240}
\definecolor{newbg}{RGB}{240,255,245}
\definecolor{boxline}{RGB}{100,116,139}

\hypersetup{
  colorlinks=true,
  linkcolor=teal!60!black,
  citecolor=teal!60!black,
  urlcolor=blue!50!black,
  pdftitle={FCS 嵌入式有序 KV 内核详细技术文档},
  pdfauthor={FCS / embdb/x}
}

\lstset{
  basicstyle=\ttfamily\footnotesize,
  keywordstyle=\bfseries\color{teal!70!black},
  commentstyle=\itshape\color{gray!55!black},
  stringstyle=\color{purple!55!black},
  breaklines=true,
  frame=single,
  rulecolor=\color{black!20},
  backgroundcolor=\color{black!3},
  columns=fullflexible,
  keepspaces=true,
  showstringspaces=false,
  aboveskip=0.6em,
  belowskip=0.5em,
  xleftmargin=0.4em
}
\lstdefinelanguage{pseudo}{
  morekeywords={function,procedure,end,if,then,else,elseif,while,for,do,return,true,false,and,or,not,in,nil,lock,unlock,atomic},
  morecomment=[l]{//},
  sensitive=false
}
\lstdefinelanguage{fcsapi}{
  morekeywords={Store,RC,TidyPolicy,Soft,Hard,Periodic,MORE,OK,open,configure_deploy,polish,set_tidy_policy,maybe_tidy,tidy_snap},
  morecomment=[l]{//},
  sensitive=true
}

\algrenewcommand\algorithmicrequire{\textbf{输入:}}
\algrenewcommand\algorithmicensure{\textbf{输出:}}
\renewcommand{\algorithmiccomment}[1]{\hfill\textit{// #1}}

\pagestyle{fancy}
\fancyhf{}
\fancyhead[L]{\small FCS 详细技术文档}
\fancyhead[R]{\small\nouppercase{\leftmark}}
\fancyfoot[C]{\thepage}
\renewcommand{\headrulewidth}{0.3pt}

% 用环境而非宏参数：lstlisting 不能进 \newcommand 参数
\newenvironment{fixold}{%
\begin{mdframed}[backgroundcolor=oldbg,linecolor=red!40!black,linewidth=0.6pt,innertopmargin=4pt,innerbottommargin=4pt]
\textbf{\color{red!60!black}修前 / 老路径}\\[0.3em]\ignorespaces
}{\end{mdframed}}
\newenvironment{fixnew}{%
\begin{mdframed}[backgroundcolor=newbg,linecolor=teal!50!black,linewidth=0.6pt,innertopmargin=4pt,innerbottommargin=4pt]
\textbf{\color{teal!50!black}修后 / 新路径}\\[0.3em]\ignorespaces
}{\end{mdframed}}

\title{\textbf{FCS：Fractal Cluster Store}\\[0.35em]
\large 嵌入式有序 KV 内核\\详细技术文档\\[0.25em]
\normalsize 代码详解 · 伪代码 · 流程图 · 踩坑对照\\[0.15em]
\footnotesize 冻结至 ROUND18 / 产品树 \texttt{embdb/x}}
\author{基于 DESIGN3 · EVAL17 · ROUND13--18 · g15--g18}
\date{2026-10-21}

\begin{document}
\maketitle

\begin{abstract}
\noindent
本文面向两类读者：\textbf{（1）}仅有基础数据库知识、需要自底向上读懂实现的工程师——见第~\ref{sec:code-deep}~章（图解 + 手算示例 + API 剧本）；\textbf{（2）}已熟悉架构、需要伪代码/流程图与踩坑修前修后对照的维护者。
配套：\texttt{fcs-pits.html}、\texttt{include/fcs.h}。标注「实测」的数字来自 ROUND/EVAL，非本编译复测。
\end{abstract}

\tableofcontents
\newpage

\section{概述与产品命题}

\subsection{定位}
FCS 是面向\textbf{固定容量持久化块设备}的嵌入式有序 KV 内核：
\begin{itemize}[leftmargin=1.6em]
  \item 键类型固定为 \texttt{u64}；值为任意字节序列（大值走溢出页链）。
  \item 公开操作：\texttt{put} / \texttt{get} / \texttt{del} / \texttt{del\_range} / 前缀扫描与游标。
  \item 热路径无堆分配倾向（内存表为连续值区 + 开放寻址）；崩溃后恢复到\textbf{某个提交点}的一致状态。
  \item \textbf{不是} B+ 树引擎，也\textbf{不是} SQL；单进程内多写者由引擎锁串行化公开入口。
\end{itemize}

\subsection{核心命题}
v2 的 B+ 树大量复杂度花在「维持平衡」的副作用上（分隔键下推、父指针、级联分裂、COW）。
FCS 的命题是：\textbf{直接维护平衡本身，于是「树」消失}。
三个同构（分形）层次：
\begin{center}
\begin{tabular}{@{}lll@{}}
\toprule
层次 & 有序结构 & 物理实体 \\
\midrule
store & 单元表 directory & cell $=$ 键区间 + 段视图 \\
run & 页表 fences & 页 $=$ 有序记录序列 \\
page & 页内记录 & 二分 / 线性 + 可选前缀压缩 \\
\bottomrule
\end{tabular}
\end{center}

关键点：\textbf{单元不拥有数据}。数据只存在于不可变的段（run）中；单元是「(run, 页区间, 键域)」的集合视图。
因此分裂/合并单元可在 $O(\#\text{段})$ 内完成，而无需自底向上传播。

\subsection{三条平衡不变量}
\begin{description}[style=nextline,leftmargin=2.2em]
  \item[B1] 单元物理字节 $\le \texttt{BMAX}$（默认 $32\,\mathrm{KiB}$）：超限则在页边界切开。
  \item[B2] 单元段数 $\le \texttt{SMAX}$：超限则压实为更少段（布隆使读对段数不敏感）。
  \item[B3] 每次提交的重写预算随垃圾产生速率自适应；候选压实单元按 0/1 背包分配。
\end{description}

\subsection{文档范围与冻结线}
\begin{itemize}[leftmargin=1.6em]
  \item \textbf{代码}：\texttt{embdb/x/include/\{fcs,common,blockdev,layout\}.h}（产品主树）。
  \item \textbf{冻结轮次}：ROUND16（P2a+P2b）$\to$ ROUND17（fsync 出锁 / hole\_fit / 家族游标）$\to$ ROUND18（护洞 / 碎片 A / 整理态；并洞 B 默认关）。
  \item \textbf{评估主文}：\texttt{docs/00-entry/EVAL17.md}；升格对照：\texttt{docs/00-entry/BATTLE.md}。
  \item \textbf{非产品默认}：本机「小预留、TRIM 关」分叉已归档，不与主树混编。
\end{itemize}

\subsection{一页结论（摘自 EVAL17）}
\begin{table}[ht]
\centering
\small
\begin{tabular}{@{}p{2.2cm}p{7.2cm}p{4.8cm}@{}}
\toprule
维度 & 结论 & 最大缺口 \\
\midrule
定位 & 追加 run + 单元目录 + 分层归并 + 分片 TRIM & 键固定 u64；单进程 \\
正确性 & 回归含 g15--g18；格式/墓碑/宽家族 & 无小时 soak；无真拔电 \\
耐久 & 双槽 SB；页复用看 \texttt{sb\_synced\_} & 阶段二失败窗口变宽 \\
写 & 整理后 $\ll$ SQLite / LMDB & 默认同步时写接近 SQLite \\
读 & 点查追平 LMDB；range 整理后仍慢于 LMDB & 未整理 range 差距大 \\
空间 & 整理后紧；合并峰 det200k $\sim 1.71\times$ & 养洞；P2c 未做 \\
延迟 & fsync 出锁；读者残余 $\sim$ 维护单元 & 锁内仍有 \texttt{flush\_mem} \\
\bottomrule
\end{tabular}
\caption{产品冻结时的能力边界（细节见第~\ref{sec:eval}~章）。}
\end{table}

\section{代码详解：从底层往上拆}
\label{sec:code-deep}

本章面向\textbf{只熟悉基础数据库概念}的读者（表、主键、事务、提交、索引即可）。
不要求读过 LSM / B+ 树论文。每一层先给「你已经知道的类比」，再落到 \texttt{embdb/x/include/fcs.h} 里的真实结构与调用链，并配图与\textbf{可手算的小例子}。

\begin{mdframed}[linecolor=boxline,linewidth=0.5pt,innertopmargin=4pt,innerbottommargin=4pt]
\textbf{阅读约定}
\begin{itemize}[leftmargin=1.4em,itemsep=0.2em]
  \item 真相源：\texttt{include/fcs.h}（模块 \texttt{M01}--\texttt{M07}）、\texttt{common.h}（\texttt{PAGE=4096}）、\texttt{blockdev.h}。
  \item 公开入口：\texttt{Store::put/get/del/commit/\ldots}；内部名出现时会标明文件位置含义。
  \item 「演示」数字是教学用的微型库，\textbf{不是} ROUND 实测 KPI。
\end{itemize}
\end{mdframed}

% =============================================================================
\subsection{先建立一张「地图」}
\label{sec:cd-map}

把 FCS 想成一张\textbf{只有一张表}的数据库：主键是 64 位整数，值是任意字节。
没有 SQL，没有多表 JOIN——但「写缓冲 → 落盘 → 查最新 → 定期整理」这条故事线，和你在 SQLite / 日志型存储里见过的一样。

\begin{figure}[H]
\centering
\begin{tikzpicture}[
  font=\small,
  L/.style={draw,rounded corners=2pt,align=center,inner sep=5pt,minimum width=2.6cm,minimum height=1.05cm},
  arr/.style={-{Latex},thick}
]
\node[L,fill=black!4] (dev) {层 0\\块设备\\连续页号};
\node[L,fill=teal!12,right=0.55cm of dev] (pg) {层 1\\数据页\\4\,KiB};
\node[L,fill=teal!18,right=0.55cm of pg] (rec) {层 2\\记录\\键+长+值};
\node[L,fill=orange!15,below=0.7cm of pg] (mem) {层 3\\MemTable\\未落盘写};
\node[L,fill=violet!12,below=0.7cm of rec] (run) {层 4\\Run\\不可变有序段};
\node[L,fill=blue!10,below=0.7cm of run] (cell) {层 5\\Cell+Seg\\目录视图};
\node[L,fill=green!12,left=0.55cm of cell] (vis) {层 6\\可见性\\版本/墓碑};
\node[L,fill=red!10,right=0.55cm of cell] (api) {层 7\\公开 API\\put/get/commit};
\draw[arr] (dev) -- (pg);
\draw[arr] (pg) -- (rec);
\draw[arr] (mem) -- (run);
\draw[arr] (run) -- (cell);
\draw[arr] (cell) -- (vis);
\draw[arr] (vis) -- (api);
\draw[arr,dashed] (mem.east) -| (api.south);
\draw[arr,dashed] (rec.south) -- (run.north);
\end{tikzpicture}
\caption{自底向上的七层拆解。虚线：写先经 MemTable；读可能先命中 MemTable 再探盘上 Run。}
\end{figure}

\begin{center}
\small
\begin{tabular}{@{}llp{7.2cm}@{}}
\toprule
你熟悉的概念 & FCS 对应物 & 一句话 \\
\midrule
数据库文件 & 一块 \texttt{BlockDev}（页号 $0,1,2,\ldots$） & 没有「文件系统里的多个 .db 文件」抽象 \\
页 / 页缓存 & $4096\,\mathrm{B}$ 页 + 两路组相联缓存 & 一切持久字节都落在某页偏移上 \\
行 & 记录（键 $8\,\mathrm{B}$ + 长度 + 值） & 键固定 \texttt{u64} \\
写缓冲 / WAL 前奏 & \texttt{MemTable} & \texttt{put} 先改内存，\texttt{commit(true)} 才变 Run \\
不可变 SST 文件 & \texttt{Run}（一段连续数据页） & 写完不再改；更新靠「更新的 Run」覆盖 \\
表分区 / 键范围 & \texttt{Cell}（$[lo,hi)$） & 目录把键空间切成区间 \\
索引指针 & \texttt{Seg}＝「某 Run 的某页区间 + 键域」 & Cell \textbf{不拥有}数据，只挂视图 \\
事务提交点 & 超级块 SB + 清单增量 & 崩溃恢复到「某个已同步的提交点」 \\
\bottomrule
\end{tabular}
\end{center}

% =============================================================================
\subsection{层 0：块设备——整库就是一串页}
\label{sec:cd-dev}

对基础读者：不要想象「一个目录里很多 .sst 文件」。FCS 面对的是\textbf{一块固定容量的块设备}（或预分配的大文件），整库地址空间就是页号 $0,1,2,\ldots$。

\begin{itemize}[leftmargin=1.4em]
  \item 页号类型 \texttt{PageId}（\texttt{u32}），页大小编译期常量 \texttt{PAGE = 4096}（\texttt{common.h}）。
  \item 读：\texttt{page(p)} 经页缓存返回指针；写：先改缓冲再按提交路径刷设备。
  \item 空闲页记在空闲表；\textbf{复用有闸}——须等对应超级块已同步（\texttt{sb\_synced\_}），否则崩溃会破坏「完整旧态」。
  \item 下文演示假设前几百页够用；超级块占双槽（交替写入），见格式章。
\end{itemize}

\begin{figure}[H]
\centering
\begin{tikzpicture}[font=\scriptsize,box/.style={draw,minimum width=1.15cm,minimum height=0.7cm,align=center,inner sep=1pt}]
\foreach \i/\lab/\c in {0/{SB-A}/red!15, 1/{SB-B}/red!10, 2/{数据}/teal!15, 3/{数据}/teal!15, 4/{...}/white, 5/{清单}/orange!15, 6/{数据}/teal!15, 7/{空闲}/black!6} {
  \node[box,fill=\c] at (\i*1.25,0) (p\i) {p\i\\\lab};
}
\draw[-{Latex}] (-0.3,-1.0) -- (9.0,-1.0) node[right] {页号增大 $\rightarrow$};
\end{tikzpicture}
\caption{设备是线性页数组。超级块两槽、数据页、清单页、空闲页混排在同一地址空间。}
\end{figure}

% =============================================================================
\subsection{层 1--2：数据页里的记录长什么样}
\label{sec:cd-rec}

\subsubsection{记录字节布局（教学用简化）}
一条小值记录在数据页上的逻辑布局（启用页内前缀前，\texttt{pfx=0}）：

\begin{center}
\begin{tikzpicture}[font=\ttfamily\footnotesize]
\node[draw,fill=black!3,inner sep=4pt] {
\begin{tabular}{@{}c@{\,}c@{\,}c@{}}
\texttt{key 8B} & \texttt{len 2B} & \texttt{value (偶数对齐)} \\
\end{tabular}
};
\end{tikzpicture}
\end{center}

\begin{itemize}[leftmargin=1.4em]
  \item 头文件：\texttt{M02}，\texttt{rec\_size(len) = 10 + 对齐后的值长}；溢出值时 \texttt{len = OVFL\_MARK (0xFFFF)}，值区换成 $16\,\mathrm{B}$ 描述符。
  \item \textbf{版本不写在每条记录里}：同一条 Run 共享同一个逻辑版本（一次 \texttt{commit} / 一次压实产物）。
  \item 页头还记 used、首末键、可选前缀压缩（有序键可少存高位字节）；旧页 \texttt{fmt=0} 走全键路径。
\end{itemize}

\subsubsection{示例 A：手算一条页上的字节}
假设写入键 $k=100$，值 \texttt{"hi"}（长度 $2$）：

\begin{enumerate}[leftmargin=1.6em]
  \item 值长 $2$，对齐后仍为 $2$；\texttt{rec\_size(2) = 10 + 2 = 12} 字节。
  \item 小端想象（仅示意）：\texttt{64 00 00 00 00 00 00 00 | 02 00 | 68 69}（键 $100$、长 $2$、ASCII \texttt{h i}）。
  \item 若同页再写 $k=101$，值 \texttt{"ok"}，再占 $12$ 字节；页从偏移 \texttt{DHDR}（数据区起点，约 $40$）起顺序追加。
\end{enumerate}

\begin{figure}[H]
\centering
\begin{tikzpicture}[font=\scriptsize]
\draw[thick] (0,0) rectangle (12,1.6);
\fill[black!8] (0,0) rectangle (1.2,1.6);
\node[align=center] at (0.6,0.8) {页头\\DHDR};
\fill[teal!25] (1.2,0) rectangle (3.0,1.6);
\node[align=center] at (2.1,0.8) {rec\\k=100};
\fill[teal!40] (3.0,0) rectangle (4.8,1.6);
\node[align=center] at (3.9,0.8) {rec\\k=101};
\fill[black!4] (4.8,0) rectangle (12,1.6);
\node at (8.4,0.8) {未用（used 之后）};
\node[anchor=west] at (0,-0.45) {一页 $=4096\,\mathrm{B}$；多条记录按键\textbf{有序}排列（刷盘时排序写入）。};
\end{tikzpicture}
\caption{示例 A：一页内两条小记录的示意位置。}
\end{figure}

% =============================================================================
\subsection{层 3：MemTable——先写内存，再谈落盘}
\label{sec:cd-mem}

\subsubsection{类比}
像「还没 \texttt{fsync} 的私有写缓冲」：\texttt{put}/\texttt{del} 立刻改内存表；同进程其它线程的 \texttt{get} 在加锁后也能看见；\textbf{掉电会丢}尚未 \texttt{commit(true)} 的小值。

实现上 MemTable 不是 \texttt{std::map}：用开放寻址 + 连续值区，避免每条记录两次堆分配（见 \texttt{M03} 注释）。对读者只需记住三点：按键覆盖写、可有序吐出、点查 $O(1)$ 期望。

\subsubsection{代码行为（\texttt{M03} / \texttt{put\_raw}）}
\begin{enumerate}[leftmargin=1.6em]
  \item 拿写锁（\texttt{WGuard}）。
  \item \texttt{mem\_.upsert(key)}：有则覆盖，无则插入。
  \item 小值：值拷进 MemTable 连续值区；大值（$>$\texttt{INLINE\_VAL\_MAX}，约 $3000$）：\textbf{先}写成溢出页链，MemTable 只留 $16\,\mathrm{B}$ 描述符。
  \item \texttt{del}：同一套 upsert，但长度记 $0$（墓碑语义，见层~6）。
  \item \textbf{此时不分配数据 Run，不改超级块。}
\end{enumerate}

\begin{lstlisting}[language=pseudo,caption={put 热路径（对照 Store::put\_raw）}]
function put_raw(key, val, len):
  lock write
  if len > INLINE_VAL_MAX:
    d = write_overflow_chain(val)     // 已占设备页
    mem.upsert(key) = descriptor(d)  // 16B，ovfl=true
  else:
    mem.upsert(key) = copy(val)
  mem[key].ver = open_vers()          // 逻辑版本
  unlock
\end{lstlisting}

\subsubsection{示例 B：三次 put，尚未 commit}
\begin{center}
\small
\begin{tabular}{@{}llll@{}}
\toprule
步骤 & 调用 & MemTable 内容（键$\to$值） & 盘上 Run \\
\midrule
1 & \texttt{put(1,"a")} & $1\to\mathrm{a}$ & （无） \\
2 & \texttt{put(2,"b")} & $1\to\mathrm{a},\;2\to\mathrm{b}$ & （无） \\
3 & \texttt{put(1,"A")} & $1\to\mathrm{A},\;2\to\mathrm{b}$（覆盖） & （无） \\
\bottomrule
\end{tabular}
\end{center}

\texttt{get(1)} 在步骤 3 之后直接命中 MemTable，返回 \texttt{"A"}，\textbf{不会}去探 Run。

% =============================================================================
\subsection{层 4：Run——落盘后的「不可变有序文件」}
\label{sec:cd-run}

\subsubsection{类比}
一条 Run $\approx$ LSM 里的一份 SSTable：键有序、写完只读、更新靠更新的文件。
物理上是连续若干数据页（可再跟元数据页存 fences / bloom 等）。

\subsubsection{结构体字段（读注释即可）}
\texttt{struct Run}（\texttt{M05}）关键字段：

\begin{center}
\small
\begin{tabular}{@{}lp{9.5cm}@{}}
\toprule
字段 & 含义 \\
\midrule
\texttt{page0, npages} & 数据页起止 \\
\texttt{first, last} & 整条 Run 键范围 \\
\texttt{ver, seq} & 逻辑版本；全局单调序号（新$\to$旧遍历用 \texttt{seq}） \\
\texttt{fences[]} & 每页首键，二分定位到页 \\
\texttt{bloom} & 「键一定不在」时可跳过整段 \\
\texttt{live\_recs / live\_bytes} & 还有多少是「某键的最新副本」（归零可整段回收） \\
\texttt{trim} & P2b：前缀页已释放，第 0 页留作身份锚 \\
\bottomrule
\end{tabular}
\end{center}

\begin{figure}[H]
\centering
\begin{tikzpicture}[font=\small,
  pg/.style={draw,rounded corners=1pt,align=center,inner sep=5pt,
             minimum width=2.0cm,minimum height=1.15cm,fill=teal!18},
  mt/.style={draw,rounded corners=1pt,align=center,inner sep=5pt,
             minimum width=2.2cm,minimum height=1.15cm,fill=orange!20}
]
% 标题单独在上方，绝不与页框重叠
\node[font=\bfseries] (ttl) at (3.5,1.85)
  {Run\quad seq$=5$\quad ver$=3$\quad page0$=10$\quad npages$=3$};
\node[pg] (p0) at (0.0,0) {p10\\首键 1};
\node[pg] (p1) at (2.4,0) {p11\\首键 50};
\node[pg] (p2) at (4.8,0) {p12\\首键 200};
\node[mt] (m)  at (7.4,0) {meta\\fences / bloom};
% 外框表示「一条物理 Run」
\draw[thick,violet!50!black,rounded corners=2pt]
  ($(p0.south west)+(-0.25,-0.35)$) rectangle ($(m.north east)+(0.25,0.35)$);
\draw[-{Latex},thick] (p0) -- (p1) -- (p2);
\draw[dashed] (p2) -- (m);
\node[anchor=west,font=\scriptsize,align=left] at (-0.3,-1.55)
  {数据页连续占用设备页号 $[10,13)$；\texttt{fences=[1,50,200]} 供二分定位；meta 可选。};
\end{tikzpicture}
\caption{一条 Run：标题在上、页框在下；数据页连续，可选 meta 挂在后面。}
\end{figure}

\paragraph{用 fences 手算「键落在哪一页」}
设 \texttt{fences = [1, 50, 200]}（每页首键）。查键 $k$ 时：找最大的 $i$ 使 $\texttt{fences}[i]\le k$，则候选页是 \texttt{page0}+i。

\begin{center}
\small
\begin{tabular}{@{}lll@{}}
\toprule
查询键 $k$ & 命中页 & 理由 \\
\midrule
$1$ & p10 & $1\le k < 50$ \\
$80$ & p11 & $50\le k < 200$ \\
$200$ & p12 & $k\ge 200$ \\
$999$（且 $>$ Run.last） & （无） & 可先被 Run 的 \texttt{[first,last]} 拒绝 \\
\bottomrule
\end{tabular}
\end{center}

Bloom 再挡一层：若 \texttt{bloom.may(k)=false}，整条 Run 都不用读页——这就是「读放大常接近 1 页」的来源。

\subsubsection{示例 C：commit 把 MemTable 刷成 Run}
接示例 B（MemTable 为 $1\to\mathrm{A},\,2\to\mathrm{b}$），调用 \texttt{commit(true)}（教学简化）：

\begin{enumerate}[leftmargin=1.6em]
  \item \textbf{锁内 phase1}：MemTable 按键有序吐出 $\to$ 写成新 Run（设 \texttt{page0=10}，\texttt{seq=1}，\texttt{ver=1}）；目录给覆盖该键域的 Cell 挂上 Seg；抓超级块载荷（尚未宣称「已同步」）。
  \item \textbf{锁外 phase2}：数据页与超级块落盘并 \texttt{sync}（详见 fsync 章）。
  \item \textbf{锁内 phase3}：发布 \texttt{sb\_synced\_}；清空已刷的 MemTable 条目。
\end{enumerate}

此后盘上逻辑内容：$(1,\mathrm{A}),\,(2,\mathrm{b})$。MemTable 变空；\texttt{get(1)} 走盘上 Run。

% =============================================================================
\subsection{层 5：Cell 与 Seg——目录，而不是「拥有数据的节点」}
\label{sec:cd-cell}

\subsubsection{为什么要多此一举？}
若每次分裂都像 B+ 树那样搬记录，代价高。
FCS 的关键句：

\begin{quote}
\textbf{单元（Cell）不拥有数据。}数据只活在 Run 里；Cell 只保存一组 \texttt{Seg}——「我关心这条 Run 的哪一段页、哪一段键」。
\end{quote}

因此切分一个 Cell 往往只需改段表（$O(\#\text{段})$），不必重写用户数据。
可以把 Cell 想成「表分区的元数据行」，把 Seg 想成「指向某个不可变文件切片的索引项」。

\subsubsection{三个结构}
\begin{description}[style=nextline,leftmargin=2.4em]
  \item[\texttt{Cell}] 键区间 $[lo,hi)$ + 有序段表 \texttt{segs} + 存活/垃圾字节账（\texttt{live}/\texttt{dead}）。
  \item[\texttt{Seg}] 字段 \texttt{(rid, p0, p1, klo, khi, live)}：指向 Run~\texttt{rid} 的页半开区间 $[p0,p1)$，且\textbf{只}对键 $[klo,khi)$ 负责。页可以和邻居重叠，键域必须切开。
  \item[\texttt{SegKey}] 排序键 $(\texttt{seq},\texttt{rid})$：段表按 \texttt{seq} 升序存放；遍历时 \texttt{rbegin()} 先拿到\textbf{最新}段。
\end{description}

\begin{figure}[H]
\centering
\begin{tikzpicture}[font=\small,
  cell/.style={draw,rounded corners=2pt,align=center,inner sep=5pt,
               minimum width=3.4cm,minimum height=1.0cm,fill=blue!12},
  seg/.style={draw,rounded corners=1pt,align=left,inner sep=4pt,fill=blue!6},
  note/.style={draw,rounded corners=2pt,align=center,inner sep=5pt,fill=violet!12}
]
% ---- 第一层：两个 Cell ----
\node[cell] (c0) at (0,3.2) {Cell$_0$\quad $[0,\,100)$};
\node[cell] (c1) at (5.6,3.2) {Cell$_1$\quad $[100,\,+\infty)$};
% ---- 第二层：各 Cell 挂的 Seg（同列对齐，互不重叠）----
\node[seg] (s0) at (0,1.6)
  {Seg $\to$ Run\,3\quad pages $[0,2)$\\键域 $k\in[0,100)$};
\node[seg] (s1a) at (5.6,1.6)
  {Seg $\to$ Run\,3\quad pages $[1,3)$\\键域 $k\in[100,300)$};
\node[seg,fill=teal!15] (s1b) at (5.6,0.2)
  {Seg $\to$ Run\,5\quad pages $[0,1)$\\键域 $k\in[100,+\infty)$\quad\textbf{更新}};
\draw[-{Latex},thick] (c0) -- (s0);
\draw[-{Latex},thick] (c1) -- (s1a);
\draw[-{Latex},thick] (c1.south) -- ++(0,-0.15) -| (s1b.north);
% ---- 第三层：物理 Run 说明（单独一行，避开上方框）----
\node[note,minimum width=9.6cm] (run) at (2.8,-1.2)
  {物理层：同一条 Run\,3 被两个 Cell 各挂一段（页可重叠，键域必须切开）};
\draw[dashed,violet!60!black] (s0.south) -- (run.north -| s0);
\draw[dashed,violet!60!black] (s1a.south) -- (run.north -| s1a);
\end{tikzpicture}
\caption{目录视图（分层绘制）：Cell $\to$ Seg $\to$ 共享物理 Run；Cell$_1$ 另有更新的 Run\,5。}
\end{figure}

\paragraph{读图要点}
\begin{itemize}[leftmargin=1.4em]
  \item Cell$_0$ 与 Cell$_1$ 都引用 Run\,3，但键域在 $100$ 切开：左边的 Seg 绝不会把键 $150$ 算进自己的 \texttt{live}。
  \item 页区间可以重叠（图中 $[0,2)$ 与 $[1,3)$ 都含页 $1$）：重叠页里同时可能有左右两边的键，\textbf{靠键域过滤}，不是靠「独占页」。
  \item Cell$_1$ 上 Run\,5 的 Seg 更新：\texttt{get} 时 \texttt{rbegin()} 先探 Run\,5；命中（含墓碑）即停，不会「救回」Run\,3 里的旧值。
\end{itemize}

\paragraph{示例：两个 get 各走哪条路}
\begin{center}
\small
\begin{tabular}{@{}llp{6.6cm}@{}}
\toprule
调用 & 进入的 Cell & 探段顺序（新$\to$旧） \\
\midrule
\texttt{get(50)} & Cell$_0$ & 仅 Run\,3 的左 Seg；在页 $[0,2)$ 内找 $50$ \\
\texttt{get(150)} & Cell$_1$ & 先 Run\,5；未命中再 Run\,3 右 Seg（页 $[1,3)$，且 $k\ge100$） \\
\bottomrule
\end{tabular}
\end{center}

\subsubsection{示例 D：一次切分在目录上发生了什么}
初始：唯一 Cell $[0,\infty)$，一段指向 Run\,3 全页。
在键 $100$ 切开后：

\begin{itemize}[leftmargin=1.4em]
  \item 左 Cell：Seg 键域收到 $[0,100)$；若某源段在左侧\textbf{零记录}，段壳必须丢掉（否则「幽灵段」，见踩坑章）。
  \item 右 Cell：对称处理 $[100,\infty)$。
  \item Run\,3 的\textbf{字节一动不动}；变的只是两份目录视图。
\end{itemize}

切分前后对照（教学）：

\begin{center}
\small
\begin{tabular}{@{}lp{5.2cm}p{5.2cm}@{}}
\toprule
 & 切分前 & 切分后 \\
\midrule
Cell 数 & 1 个 $[0,\infty)$ & 2 个 $[0,100)$ / $[100,\infty)$ \\
Run\,3 数据页 & 原样 & 原样（无重写） \\
Seg 键域 & $[0,\infty)$ & 左 $[0,100)$，右 $[100,\infty)$ \\
\bottomrule
\end{tabular}
\end{center}

% =============================================================================
\subsection{层 6：可见性——版本、墓碑、范围删除}
\label{sec:cd-vis}

\subsubsection{点查规则（你需要记住的唯一公式）}
对键 $k$、视图版本 $\mathit{view}$（通常 $=$ 当前打开版本）：

\begin{enumerate}[leftmargin=1.6em]
  \item 若 MemTable 有 $k$ 且其 \texttt{ver} $\le \mathit{view}$：以内存为准（len$=0$ $\Rightarrow$ 不存在）。
  \item 否则找到含 $k$ 的 Cell，按段从新到旧（\texttt{segs.rbegin}）：
  \begin{itemize}
    \item bloom 说「不可能」$\to$ 跳过；
    \item \texttt{probe} 在段页区间内找到记录：len$=0$ 是墓碑 $\to$ \texttt{NOTFOUND}；否则返回值；
    \item 找不到 $\to$ 试更老的段。
  \end{itemize}
  \item 都没有 $\to$ \texttt{NOTFOUND}。
\end{enumerate}

\begin{lstlisting}[language=pseudo,caption={get\_at 主循环（对照 Store::get\_at）}]
function get_at(key, view):
  if mem has key and mem.ver <= view:
    return mem value or NOTFOUND if tombstone
  C = find_cell(key)
  for seg in C.segs from newest to oldest:
    R = run(seg.rid)
    if not bloom.may(key): continue
    rc, val = probe(R, seg.pages, key, view)
    if found tombstone: return NOTFOUND
    if found value: return value
  return NOTFOUND
\end{lstlisting}

\subsubsection{示例 E：覆盖与删除的时间线}
\begin{center}
\small
\begin{tabular}{@{}clp{8.6cm}@{}}
\toprule
时刻 & 操作 & 读者看见的 $k=7$ \\
\midrule
T0 & （空库） & 不存在 \\
T1 & \texttt{put(7,"x"); commit} & \texttt{"x"}（Run$_1$） \\
T2 & \texttt{put(7,"y"); commit} & \texttt{"y"}（Run$_2$ 更新；Run$_1$ 上旧值变垃圾） \\
T3 & \texttt{del(7); commit} & 不存在（Run$_3$ 墓碑；probe 遇 len$=0$ 即停） \\
\bottomrule
\end{tabular}
\end{center}

\begin{figure}[H]
\centering
\begin{tikzpicture}[font=\scriptsize,r/.style={draw,minimum width=2.4cm,minimum height=0.85cm,align=center}]
\node[r,fill=teal!20] (a) {Run$_1$\\{7 $\rightarrow$ x}};
\node[r,fill=teal!35,right=0.4cm of a] (b) {Run$_2$\\{7 $\rightarrow$ y}};
\node[r,fill=red!20,right=0.4cm of b] (c) {Run$_3$\\{7 $\rightarrow$ 墓碑}};
\draw[-{Latex}] (a) -- (b) -- (c);
\node[anchor=west] at (0,-1.1) {get 从右（新）往左探：先碰到墓碑 $\Rightarrow$ 结束，不会返回 x/y。};
\end{tikzpicture}
\caption{示例 E：同一键多版本并存；读只认最新段上的结果。}
\end{figure}

范围删除（\texttt{del\_range}）还会在清单里留下 \texttt{M\_RANGE} 类记录：不占用户键空间，但读路径要用 \texttt{range\_hidden} 判断键是否被区间墓碑盖住。

% =============================================================================
\subsection{层 7：把 API 串成一条「演示剧本」}
\label{sec:cd-demo}

下面用\textbf{完整微型剧本}把层 0--6 串起来。假设教学库键很小，一切发生在一个 Cell 内。
每一步给出「MemTable / 盘上 Run / get 结果」三列状态，便于对照代码。

\subsubsection{剧本步骤与状态表}
\begin{enumerate}[leftmargin=1.6em]
  \item \texttt{open(dev)}：读双槽超级块，重放清单，重建 \texttt{runs\_} 与 Cell 目录；空库则一个 Cell 覆盖全域。
  \item \texttt{put(10, "cat")}；\texttt{put(20, "dog")} —— 仅 MemTable。
  \item \texttt{commit(true)} —— 产生 Run\,A（两键），SB 持久；MemTable 清空。
  \item \texttt{put(10, "CAT")}；\texttt{del(20)} —— 又进 MemTable。
  \item \texttt{get(10)} $\to$ \texttt{"CAT"}（内存命中）；\texttt{get(20)} $\to$ 不存在（内存墓碑）。
  \item \texttt{commit(true)} —— 产生 Run\,B；目录两段，新段在上。
  \item 再次 \texttt{get(10)}：MemTable 空，探 Cell 最新段 Run\,B 得 \texttt{"CAT"}；\texttt{get(20)} 在 Run\,B 见墓碑 $\to$ 不存在。
\end{enumerate}

\begin{center}
\small
\begin{tabular}{@{}clp{3.4cm}p{3.6cm}@{}}
\toprule
步 & MemTable & 盘上 Run（新$\to$旧） & \texttt{get(10)} / \texttt{get(20)} \\
\midrule
1 & $\emptyset$ & （无） & 不存在 / 不存在 \\
2 & $10\to\mathrm{cat},\,20\to\mathrm{dog}$ & （无） & cat / dog \\
3 & $\emptyset$ & A: $10\to\mathrm{cat},20\to\mathrm{dog}$ & cat / dog \\
4 & $10\to\mathrm{CAT},\,20\to$墓碑 & A（同上） & CAT / 不存在 \\
5 & （同 4） & A & CAT / 不存在 \\
6--7 & $\emptyset$ & B（新）, A（旧） & CAT / 不存在 \\
\bottomrule
\end{tabular}
\end{center}

\begin{figure}[H]
\centering
\begin{tikzpicture}[
  font=\scriptsize,
  node distance=0.5cm and 0.55cm,
  st/.style={draw,rounded corners=2pt,align=center,inner sep=4pt,minimum width=2.2cm,fill=black!3},
  go/.style={-{Latex}}
]
\node[st] (o) {open};
\node[st,right=of o] (p1) {put 10,20};
\node[st,right=of p1] (c1) {commit\\Run A};
\node[st,below=of c1] (p2) {put/del\\Mem};
\node[st,left=of p2] (g) {get 命中\\Mem};
\node[st,left=of g] (c2) {commit\\Run B};
\node[st,below=of c2] (g2) {get 探盘\\Run B};
\draw[go] (o) -- (p1) -- (c1) -- (p2) -- (g) -- (c2) -- (g2);
\end{tikzpicture}
\caption{演示剧本控制流（省略锁与维护）。}
\end{figure}

\paragraph{步 7 的 probe 内心独白（教学）}
\begin{lstlisting}[language=pseudo,caption={get(10) 在步 7：Mem 已空}]
C = find_cell(10)                 // 唯一 Cell
for seg in C.segs from newest:    // 先 Run B，再 Run A
  if bloom of Run B may(10):
    page = fences_locate(Run B, 10)
    rec = scan_page(page, 10)
    if rec found: return "CAT"    // 停；绝不看 Run A 的旧 "cat"
\end{lstlisting}

\subsubsection{对应代码入口速查}
\begin{center}
\small
\begin{tabular}{@{}llp{6.8cm}@{}}
\toprule
调用 & 大致位置（\texttt{fcs.h}） & 你要盯的不变量 \\
\midrule
\texttt{put}/\texttt{del} & \texttt{M07h} & 只动 MemTable（大值可先写溢出链） \\
\texttt{get} & \texttt{M07e} \texttt{get\_at} & 先 mem 后段新$\to$旧；墓碑停 \\
\texttt{commit} & \texttt{commit\_phase1/2/3} & 逻辑可见 vs 耐久；锁外 sync \\
\texttt{flush\_mem} & 写路径内部 & Mem $\to$ 一条 L0 Run \\
\texttt{find\_cell} & 目录 & 键落在哪个 $[lo,hi)$ \\
\texttt{probe} & 段内定位 & fences/sp + 页内查找 \\
\bottomrule
\end{tabular}
\end{center}

\begin{lstlisting}[language=fcsapi,caption={读者可粘贴的最小用法骨架}]
Store s;
s.open(dev);
s.put(10, "cat", 3);
s.put(20, "dog", 3);
s.commit(true);                 // 耐久提交
s.put(10, "CAT", 3);
s.del(20);
char buf[32]; u16 n = 0;
s.get(10, buf, sizeof(buf), &n); // n=3, buf="CAT"
s.commit(true);
s.polish();                     // 可选：对外服务前整理
\end{lstlisting}

% =============================================================================
\subsection{提交点、清单与崩溃：和「事务」对齐的那一层}
\label{sec:cd-commit}

\begin{itemize}[leftmargin=1.4em]
  \item \texttt{commit(false)}：只推进\textbf{逻辑}提交（组可见性），可不做耐久 IO。
  \item \texttt{commit(true)}：物理提交——刷 MemTable（若需要）+ 写清单增量 + 写超级块 + 同步。
  \item 崩溃后：挂载选中合法 SB 槽，重放清单，恢复到\textbf{最后一次已同步提交}；其后的逻辑提交丢失——与「同步语义」一致。
  \item ROUND17 起：页复用必须等待 \texttt{sb\_synced\_}，避免「在飞 SB」时把旧态页抢去写坏。
\end{itemize}

\begin{figure}[H]
\centering
\begin{tikzpicture}[font=\scriptsize,box/.style={draw,align=center,inner sep=4pt,rounded corners=1pt}]
\node[box,fill=orange!15] (p1) {phase1\\{锁内：内存+发写+抓 SB}};
\node[box,fill=violet!15,right=0.6cm of p1] (p2) {phase2\\{锁外：等 IO + sync SB}};
\node[box,fill=teal!15,right=0.6cm of p2] (p3) {phase3\\{锁内：发布 synced}};
\draw[-{Latex}] (p1) -- (p2) -- (p3);
\node[anchor=west,align=left] at (0,-1.15) {读者在 phase2 期间不再被「双 fsync」挡住（见 R17）；\\get 仍可能与维护单元短暂互斥。};
\end{tikzpicture}
\caption{耐久 commit 的三阶段（教学级；细节见 fsync 专章）。}
\end{figure}

% =============================================================================
\subsection{维护与合并：为什么库还会「自己写盘」}
\label{sec:cd-maint}

基础读者常困惑：我没调用 \texttt{put}，设备灯为什么还在闪？

\begin{itemize}[leftmargin=1.4em]
  \item 覆盖写与删除会让旧 Run 变成垃圾；段数变多会让读变慢。
  \item 引擎在提交后按策略做\textbf{维护}：把多个旧 Run 归并成更少、更新的 Run（分层 / 压实），并更新目录。
  \item \texttt{configure\_deploy()} / Soft 整理态：常驻轻量整理；\texttt{polish()}：发布前重锤。
  \item 养洞（hole\_fit、碎片感知切片等）影响「新页落在文件中部还是尾部」，从而影响峰值空间——这是空间议题，不改变 get/put 语义。
\end{itemize}

\begin{figure}[H]
\centering
\begin{tikzpicture}[font=\scriptsize,box/.style={draw,inner sep=3pt,align=center,rounded corners=1pt}]
\node[box,fill=violet!12] (r1) {Run 旧\\很多垃圾};
\node[box,fill=violet!12,right=0.25cm of r1] (r2) {Run 中};
\node[box,fill=violet!12,right=0.25cm of r2] (r3) {Run 新};
\node[box,fill=green!15,right=1.0cm of r3] (out) {归并产物\\更少页};
\draw[-{Latex}] (r3.east) -- (out.west) node[midway,above] {merge};
\node[anchor=north,align=left] at ($(r2.south)!0.5!(out.south)+(0,-0.35)$) {目录改挂新 Seg；旧 Run live 归零后释放页。};
\end{tikzpicture}
\caption{维护不改变「最新键值」语义，只改物理布局与垃圾比例。}
\end{figure}

% =============================================================================
\subsection{源码目录：按模块找代码}
\label{sec:cd-modules}

\begin{center}
\small
\begin{tabular}{@{}clp{8.4cm}@{}}
\toprule
模块 & 主题 & 建议阅读顺序 \\
\midrule
\texttt{M01} & 容量/层/特性宏 & 先扫常量：\texttt{BMAX}/\texttt{SMAX}/\texttt{PAGE} 相关 \\
\texttt{M02} & 记录编解码 & 与本节层 1--2 对照 \\
\texttt{M03} & MemTable & 与 \texttt{put\_raw} 对照 \\
\texttt{M04} & Bloom & 理解 get 为何常「探一段就停」 \\
\texttt{M05} & Run/Seg/Cell & 本节层 4--5 的结构体注释极有价值 \\
\texttt{M06} & XLock / TLS & 多读者时再读；单线程可跳过 \\
\texttt{M07a--} & Store 实现 & 按 API 跳读：读$\to$读$\to$commit$\to$挂载 \\
\bottomrule
\end{tabular}
\end{center}

辅助头文件：
\begin{itemize}[leftmargin=1.4em]
  \item \texttt{common.h}：\texttt{RC}、\texttt{PAGE}、CRC/哈希。
  \item \texttt{blockdev.h}：设备读写抽象。
  \item \texttt{layout.h}：更偏 XDB 叶页槽布局；FCS 数据页以 \texttt{fcs.h} 的数据页头为准，读格式章即可。
\end{itemize}

% =============================================================================
\subsection{一张「点查」全路径图（汇总）}
\label{sec:cd-getfig}

\begin{figure}[H]
\centering
\begin{tikzpicture}[
  font=\scriptsize,
  d/.style={draw,align=center,inner sep=3pt,rounded corners=1pt,minimum width=2.3cm},
  styY/.style={fill=green!12},
  styN/.style={fill=red!10},
  styQ/.style={fill=yellow!15}
]
\node[d,styQ] (start) {get(key)};
\node[d,below=0.4cm of start] (mem) {MemTable\\有 key?};
\node[d,styY,right=1.1cm of mem] (memhit) {返回 mem\\或墓碑};
\node[d,below=0.45cm of mem] (cell) {find\_cell};
\node[d,below=0.45cm of cell] (loop) {段：新到旧};
\node[d,styQ,right=1.1cm of loop] (bloom) {bloom.may?};
\node[d,below=0.45cm of loop] (probe) {probe 页内};
\node[d,styY,right=1.1cm of probe] (ok) {返回值};
\node[d,styN,left=1.1cm of probe] (tomb) {墓碑\\NOTFOUND};
\node[d,styN,below=0.45cm of probe] (miss) {NOTFOUND};
\draw[-{Latex}] (start) -- (mem);
\draw[-{Latex}] (mem) -- node[above,font=\tiny]{是} (memhit);
\draw[-{Latex}] (mem) -- (cell);
\draw[-{Latex}] (cell) -- (loop);
\draw[-{Latex}] (loop) -- (bloom);
\draw[-{Latex}] (bloom) -- node[right,font=\tiny]{可能} (probe);
\draw[-{Latex}] (probe) -- (ok);
\draw[-{Latex}] (probe) -- (tomb);
\draw[-{Latex}] (probe) -- (miss);
\draw[-{Latex}] (loop.east) -- ++(0.35,0) |- (bloom.east);
\end{tikzpicture}
\caption{点查决策图（与代码分支一致的教学版）。}
\end{figure}

% =============================================================================
\subsection{一张「写+提交」全路径图（汇总）}
\label{sec:cd-putfig}

\begin{figure}[H]
\centering
\begin{tikzpicture}[
  font=\scriptsize,
  d/.style={draw,align=center,inner sep=3pt,rounded corners=1pt}
]
\node[d,fill=orange!15] (put) {put/del};
\node[d,right=0.5cm of put] (mt) {更新 MemTable};
\node[d,right=0.5cm of mt] (c) {commit?};
\node[d,fill=violet!12,below=0.55cm of c] (flush) {flush\_mem\\写成 Run};
\node[d,fill=violet!12,right=0.5cm of flush] (man) {清单增量};
\node[d,fill=violet!12,right=0.5cm of man] (sb) {写 SB + sync};
\node[d,fill=green!12,below=0.55cm of flush] (dir) {目录挂 Seg\\活账更新};
\draw[-{Latex}] (put) -- (mt) -- (c);
\draw[-{Latex}] (c) -- node[right,font=\tiny]{durable} (flush);
\draw[-{Latex}] (flush) -- (man) -- (sb);
\draw[-{Latex}] (flush) -- (dir);
\end{tikzpicture}
\caption{写路径教学图：未 commit 前世界只在 MemTable（大值链除外）。}
\end{figure}

% =============================================================================
\subsection{自测题（建议不看答案先做）}
\label{sec:cd-quiz}

\begin{enumerate}[leftmargin=1.6em]
  \item 只 \texttt{put} 不 \texttt{commit}，掉电后数据还在吗？大值溢出链呢？
  \item 为什么 Cell 切分通常不必重写 Run 数据页？
  \item \texttt{get} 在某一段上看到墓碑，会不会继续看更老的段把旧值「救回」？
  \item 同一物理页能否出现在两个 Cell 的 Seg 里？键域要满足什么？
\end{enumerate}

\paragraph{简答}
\begin{enumerate}[leftmargin=1.6em]
  \item 普通小值：不在。大值链若已写出但未进成功提交，恢复时不应变成用户可见数据；实现上靠提交点与释放闸，细节见耐久章。
  \item 因为数据在 Run，Cell 只改 Seg 视图。
  \item 不会——墓碑表示「更新版本删除」，必须停。
  \item 能；各自 \texttt{[klo,khi)} 必须切开，且只解释自己键域内的记录。
\end{enumerate}

% =============================================================================
\subsection{常见误解（读代码前先澄清）}
\label{sec:cd-myths}

\begin{center}
\small
\begin{tabular}{@{}p{4.2cm}p{7.6cm}@{}}
\toprule
误解 & 实际情况 \\
\midrule
「Cell 像 B+ 树节点，分裂要搬数据」 & Cell 是目录视图；分裂改 Seg，数据页可不动 \\
「Run 更新会改旧页」 & Run 不可变；更新 $=$ 新 Run + 目录挂新 Seg \\
「删除会立刻腾出设备页」 & 先留墓碑；等 live 归零 / 归并后才释放 \\
「commit 一定很慢」 & \texttt{commit(false)} 可只做逻辑提交；耐久才 sync \\
「多段一定很慢」 & Bloom + fences 让多数点查仍接近一页 \\
\bottomrule
\end{tabular}
\end{center}

\subsection{下一章怎么读}
读完本章后：
\begin{itemize}[leftmargin=1.4em]
  \item 用架构章对照「三层分形」术语是否对得上；
  \item 用格式章看 SB/清单字节；
  \item 用伪代码/流程图章看合并与 fsync；
  \item 用踩坑章理解「目录切分 / 可见性 / 提交点」写错会长什么样。
\end{itemize}
\section{体系结构}

\subsection{逻辑视图}
\begin{figure}[ht]
\centering
\begin{tikzpicture}[
  font=\small,
  box/.style={draw=black!50,rounded corners=2pt,align=center,inner sep=6pt,minimum width=2.8cm},
  arr/.style={-{Latex},thick,black!60}
]
\node[box,fill=teal!8] (store) {Store\\公开 API / 锁 / 维护};
\node[box,below=0.55cm of store,fill=cyan!6] (dir) {Cell 目录\\键区间二分};
\node[box,below=0.55cm of dir,fill=orange!6] (seg) {Seg 视图\\(run, 页区间, 键域)};
\node[box,below left=0.7cm and -0.2cm of seg,fill=black!4] (run) {Run\\fences / Bloom / sp};
\node[box,below right=0.7cm and -0.2cm of seg,fill=black!4] (page) {Page\\记录 / 前缀 / CRC};
\node[box,below=1.5cm of seg,fill=violet!6] (man) {Manifest + 双槽 SB};
\draw[arr] (store) -- (dir);
\draw[arr] (dir) -- (seg);
\draw[arr] (seg) -- (run);
\draw[arr] (seg) -- (page);
\draw[arr] (store) -- (man);
\end{tikzpicture}
\caption{FCS 逻辑层次：目录与清单是控制面，run/page 是数据面。}
\end{figure}

\subsection{组件职责}
\begin{description}[style=nextline,leftmargin=1.4em]
  \item[MemTable] 提交前的写缓冲：点查、覆盖写、有序吐出、整体清空；满则 Flush 成 L0 run。
  \item[Run] 一次提交或一次压实/合并的产物；不可变；带稀疏索引 \texttt{sp}、可选元数据页、布隆。
  \item[Cell / Seg] 逻辑可见性窗口；段键用 $(\mathrm{seq},\mathrm{rid})$ 等全序标识，避免页号复用撞身份。
  \item[Version / Manifest] 有效 run 集合与下一文件号；\textbf{恢复真相源}。每次 Flush/Compact 写新清单，再 fsync，再删过期页。
  \item[BlockDev] RamDev / FileDev / NorSim；写放大计量、故障注入、\texttt{sync}/\texttt{TRIM}。
\end{description}

\subsection{并发模型（P1）}
\begin{itemize}[leftmargin=1.6em]
  \item 公开入口默认经 \texttt{XLock}：写可重入独占；读可共享（\texttt{set\_shared\_reads}）。
  \item 写者优先；glibc \texttt{shared\_mutex} 的读者优先曾导致写者饥饿（坑~86），故自研等待队列。
  \item 每线程 TLS 槽（风险指针、共享深度、私有页缓冲）；槽耗尽退化为独占。
  \item mode~3：后台维护线程；与读者让路、预算/时间片协同（ROUND13/14）。
\end{itemize}

\subsection{维护模式}
\begin{center}
\begin{tabular}{@{}clp{8.5cm}@{}}
\toprule
模式 & 名称 & 行为摘要 \\
\midrule
0 & 提交前 & 老路径；部署契约避免退回 \\
1 & 提交后 & 默认；墙钟预算（如 $20\,\mathrm{ms}$）内做有界单元 \\
3 & 后台线程 & 缩短 commit 墙钟；读者 p99 可能变差 \\
\bottomrule
\end{tabular}
\end{center}
分步 API：\texttt{consolidate\_step} / \texttt{compact\_file\_step} 返回 \texttt{MORE}，每步 $=$ 提交点。

\section{持久化格式（FCS2）}

\subsection{从 FCSB 到 FCS2}
第~16~轮完成 \textbf{P2 格式 v2}（设计稿 \texttt{P2\_DESIGN.md}）：
\begin{itemize}[leftmargin=1.6em]
  \item 超级块魔数 \texttt{'FCS2'}，头扩至约 $48\,\mathrm{B}$ 量级字段布局（实现以头文件为准）。
  \item 清单项 $48\,\mathrm{B}$：\texttt{seq}/\texttt{ver} 升为 \textbf{64~位}，消除坑~102/103 类回绕静默错误的格式根因。
  \item 老库挂载可只读打开 v1，写路径触发升级；golden 镜像分 \texttt{tests/golden}（v2）与 \texttt{golden/v1}。
\end{itemize}

\subsection{超级块（双槽）}
\begin{itemize}[leftmargin=1.6em]
  \item 两槽交替；CRC 覆盖头 + 内联增量。
  \item 提交点语义：有效槽的状态字单次对齐写；挂载择优。
  \item ROUND17 起区分代数：\texttt{sb\_built\_}（锁内抓载荷）与 \texttt{sb\_synced\_}（锁外已持久）；页复用闸比较后者。
\end{itemize}

\subsection{清单操作码（节选）}
\begin{center}
\begin{tabular}{@{}clp{8.2cm}@{}}
\toprule
码 & 名 & 含义 \\
\midrule
1--5 & ADD/DEL/\ldots & run / split / seg 增减 \\
6 & \texttt{M\_RANGE} & 范围墓碑（只活在清单，不占键空间） \\
7 & \texttt{M\_RUN\_TRIM} & P2b：源 run 前缀释放；第~0~页为身份锚 \\
8 & \texttt{M\_RUN\_HOLE} & P2c 预留：run 内部释页（本冻结轮\textbf{未实现}） \\
\bottomrule
\end{tabular}
\end{center}

\subsection{页与记录}
\begin{itemize}[leftmargin=1.6em]
  \item 页大小 \texttt{PAGE}（与 XDB 共用布局常数）；数据页含 used、首末键、可选页内前缀（\texttt{pfx}/\texttt{fmt}）。
  \item 记录头瘦身：键 $8\,\mathrm{B}$ + 长度 $2\,\mathrm{B}$；版本是 run 属性；页 CRC 覆盖载荷。
  \item 大值：\texttt{len = OVFL\_MARK}，体为 $16\,\mathrm{B}$ 描述符，值在溢出链；\texttt{put} 时可先落链。
\end{itemize}

\subsection{身份与复用}
页号会复用，故段身份必须带 \texttt{seq}（及 rid）。提交失败前释放的页进 \texttt{pending\_free\_}，须待对应 SB 代数持久后才可复用——否则崩溃会破坏「完整旧态」。

\section{写路径、读路径与空间}

\subsection{写路径}
\begin{enumerate}[leftmargin=1.8em]
  \item \texttt{put}/\texttt{del}/\texttt{del\_range} 进入 MemTable（及待提交范围墓碑）。
  \item \texttt{commit}：\textbf{阶段一}（持锁）\texttt{flush\_mem} 编码并投递写、更新清单内存视图、抓 SB 载荷并发票号。
  \item \textbf{阶段二}（\textbf{锁外}，按票号串行）：\texttt{wait\_io} $\to$ \texttt{sync} $\to$ 写 SB $\to$ \texttt{sync}。
  \item \textbf{阶段三}（持锁）：发布 \texttt{sb\_synced\_}/\texttt{durable\_vers\_}，\texttt{promote\_pending}，可选自动归一。
\end{enumerate}
语义不变：\texttt{commit} 返回 $\Rightarrow$ 已持久；读者可见性本就可早于持久（未提交 mem 可读）。
变宽的是阶段二失败窗口：内存可能已切新视图而盘仍旧点 $\Rightarrow$ \texttt{dirty\_} + 按盘重建。

异步写：后台 IO 线程写页；同页顺序由入队保证；设备须支持异页并发。

\subsection{读路径}
\begin{enumerate}[leftmargin=1.8em]
  \item 目录二分定位 cell。
  \item 段按版本/全序从新到旧探测；Bloom 过滤负查。
  \item fence / sp 定位页；解码记录（含前缀还原）；溢出则跟链。
  \item 范围墓碑与点墓碑参与可见性；快照视图用 \texttt{view} 过滤。
\end{enumerate}
家族虚拟源（ROUND17）：同 \texttt{seq} 且键域不相交的片串成链，一链一游标顺序读，避免源 $>64$ 时掩码溢出导致 live 账漂移。

\subsection{空间与高水位}
\begin{itemize}[leftmargin=1.6em]
  \item 空闲表内存重建；快照 pin 的页不回收。
  \item 用户写入预留约 $1/8$ 容量给维护（历史规矩）。
  \item 高水位回收：ROUND15 \textbf{二步搬家}——规划与「一提交点搬一条」分离；快照钉住记入 \texttt{hw\_blockers}。
  \item \texttt{shrink} 必须落在提交点之后（坑~104：否则掉电可丢已提交数据）。
\end{itemize}

\subsection{范围墓碑要点}
墓碑是清单项，不进键空间。同 \texttt{lo} 多条须按 $(\mathrm{lo},\mathrm{ver})$ 索引，避免窄墓碑顶掉宽墓碑。
P2b 同 \texttt{seq} 家族片打破「组内 seq 连续 $\Rightarrow$ 更老」假设，探测改为 \texttt{older\_may\_contain\_grp}（坑~112 / g16 tomb）。

\section{分层合并与 P2b 养洞}
\label{sec:merge}

\subsection{层规则（摘要）}
层合并在「读放大 / 垃圾 / 悬崖」等条件下触发（ROUND13/17 有 (a)(c)(d)(e) 与 $4t$ 悬崖等旋钮）。
ROUND17 对「拧时机旋钮压峰」给出\textbf{否定结论}：简单改悬崖/脏度门槛在 layout 上往往更差；
真正能动峰值的是\textbf{分配侧}——中等产物是否进洞、片是否适配碎片洞。

\subsection{P2b：分片 + \texttt{M\_RUN\_TRIM}}
问题：整组合并先完整写出再释放源 $\Rightarrow$ 峰值接近 $2\times$ 活数据，且洞来得太晚。

做法：
\begin{enumerate}[leftmargin=1.8em]
  \item \texttt{choose\_part} 按可用洞挑选有界键区间「片」。
  \item 片封口后写 \texttt{M\_RUN\_TRIM}：推进源 run 的逻辑起点；\textbf{第~0~页永留作身份锚}。
  \item 跨重开重放 TRIM；abort/rollback/快照作废规则由 g16 覆盖。
\end{enumerate}

代价：run/片数量上升；家族同 seq 语义更复杂（见上章墓碑）。

\subsection{hole\_fit（ROUND17）}
旧逻辑仅对「大产物」（$\ge 1024$ 页）按洞分片；\textbf{中等产物}常直接落尾，错过合并刚释放的大洞。
\texttt{set\_part\_hole\_fit(true)}（部署 Soft 默认开）让中等产物也可按洞切。

\begin{table}[ht]
\centering
\small
\begin{tabular}{@{}lcc@{}}
\toprule
场景 & 大约峰值（活倍数） & 备注 \\
\midrule
修前 layout 200k & $1.96$--$2.00\times$ & 归档 \\
hole\_fit 后 & $1.75$--$1.81\times$ & 零写放大量级 \\
800k & $1.40\to\sim1.36\times$ & 同轮 \\
\bottomrule
\end{tabular}
\caption{hole\_fit 对 layout 峰值的影响（ROUND17 实测摘要）。}
\end{table}

\subsection{碎片感知切片 A（ROUND18）}
当最大洞 $< \texttt{part\_lo}$，但空闲总量足够且最大洞 $\ge \texttt{frag\_floor}$（默认 16）时，
临时在 $[\texttt{frag\_floor},\texttt{part\_lo})$ 按洞挑片（\texttt{part\_frags++}）。

\begin{table}[ht]
\centering
\small
\begin{tabular}{@{}lccc@{}}
\toprule
配置 & 峰值$\times$ & \texttt{part\_frags} & 写 \\
\midrule
frag=0 det200k & $1.82$ & 0 & $41.8\,\mathrm{MB}$ \\
frag=1 det200k & $\mathbf{1.71}$ & 2 & $41.8\,\mathrm{MB}$ \\
\bottomrule
\end{tabular}
\caption{碎片 A：$-0.11\times$，写字节齐平（ROUND18）。}
\end{table}

\subsection{有界并洞 B（ROUND18，默认关）}
先搬挡路 run 再合并空洞；机制正确（g18 \texttt{t\_coal}），但 det200k 峰 $1.71\to1.74$（写放大抬峰）。
\textbf{产品默认关闭}；Soft/\texttt{configure\_deploy} 不启用。需要时显式 \texttt{set\_part\_coalesce(true)}。

\subsection{护洞（ROUND18）}
用户路径 \texttt{alloc\_pages} 跳过 $\ge\texttt{hole\_protect\_pages}$（默认 64）的中部大洞；
维护/流式仍可进洞；NOSPACE 回退。

A/B 表明：skips 计数生效，但\textbf{200k 峰值几乎不动}（峰主因仍是养洞长尾等）；
800k 约 $-0.04\times$。故护洞作 Soft 默认（零写放大、机制正确），\textbf{不当作 200k KPI}。

\section{fsync 出锁与整理状态机}
\label{sec:fsync-tidy}

\subsection{fsync 出锁（ROUND17 / DISCUSS15）}
\subsubsection{动机}
提交者在 XLock 内做两次 \texttt{sync} 时，共享读被挡住的时间 $\approx$ 设备 fsync 墙钟。
慢盘注入（sync 睡 $40\,\mathrm{ms}\times 2$）上，读者 \texttt{get} 平均约 $79\,\mathrm{ms}$；出锁后约 $0.01\,\mathrm{ms}$。

\subsubsection{三段式提交}
\begin{center}
\begin{tabular}{@{}clp{7.8cm}@{}}
\toprule
阶段 & 锁 & 工作 \\
\midrule
phase1 & 内 & 就绪检查、\texttt{flush\_mem}、清单更新、抓 SB 载荷、发票号 \\
phase2 & \textbf{外} & 按票号串行：\texttt{wait\_io} / sync / 写 SB / sync；失败粘性 \\
phase3 & 内 & 发布 \texttt{sb\_synced\_}、\texttt{durable\_vers\_}、\texttt{promote\_pending} \\
\bottomrule
\end{tabular}
\end{center}

持锁调用方（维护提交、abort、snapshot 等）走 \texttt{commit\_locked} 内联三段，不放锁。
\texttt{durable\_vers()} 为跨线程持久进度的权威接口。

\subsubsection{已修坑（节选）}
阶段二失败后空提交短路、\texttt{cfail\_} 粘性、空票发布版本、\texttt{open}/\texttt{reload} 与顺序器交互等——均有 g17 段覆盖。
\textbf{未做}：组提交流水线、双 MemTable（锁内 \texttt{flush\_mem} 仍在）。

\subsection{整理状态机（ROUND18）}
部署侧需要「日常读友好」与「对外服务前重锤」的一等公民策略，避免漏开 hole\_fit/护洞等旋钮。

\begin{description}[style=nextline,leftmargin=1.6em]
  \item[\texttt{TidyPolicy::Off}] 不自动整理。
  \item[\texttt{Soft}] 调用 \texttt{configure\_deploy}：async + cheap 合并 + hole\_fit + 碎片 A + 护洞 + 可选 hw\_reclaim；组提交仍由应用拼。
  \item[\texttt{Hard}] 压力下 \texttt{maybe\_tidy}$\to$\texttt{consolidate\_step}。
  \item[\texttt{Periodic}] \texttt{maybe\_tidy}$\to$\texttt{compact\_file\_step}。
  \item[\texttt{polish()}] \texttt{consolidate(0)} + \texttt{shrink()}：发布前重锤。
\end{description}

观测：\texttt{tidy\_snap()} 返回 policy、runs、used/live、layout 百分比、阈值、\texttt{protect} 等。

\begin{lstlisting}[language=fcsapi,caption={部署侧最小用法示意}]
Store s;
s.open(dev);
s.configure_deploy();              // Soft 常驻
// ... 业务写入与组提交 ...
s.polish();                        // 对外读服务前
// 或：s.set_tidy_policy(TidyPolicy::Periodic);
// while (s.maybe_tidy() == RC::MORE) { }
\end{lstlisting}

\section{流程图、设计图与核心伪代码}
\label{sec:diagrams}

\subsection{端到端写路径（部署视角）}
\begin{figure}[H]
\centering
\begin{tikzpicture}[
  font=\small,
  node distance=0.55cm and 0.9cm,
  box/.style={draw=boxline,rounded corners=2pt,align=center,inner sep=5pt,minimum width=2.4cm,minimum height=0.85cm},
  io/.style={box,fill=cyan!12},
  cpu/.style={box,fill=orange!10},
  disk/.style={box,fill=violet!10},
  dec/.style={diamond,draw=boxline,aspect=2.2,inner sep=1pt,align=center,fill=amber!8},
  arr/.style={-{Latex},thick,black!55}
]
\node[io] (put) {put/del\\MemTable};
\node[cpu,right=of put] (c1) {phase1\\持锁};
\node[disk,right=of c1] (c2) {phase2\\锁外 sync};
\node[cpu,right=of c2] (c3) {phase3\\发布};
\node[io,below=1.1cm of c1] (maint) {maint\\预算/步进};
\node[disk,below=1.1cm of c3] (trim) {TRIM/洞\\choose\_part};
\draw[arr] (put) -- (c1);
\draw[arr] (c1) -- (c2);
\draw[arr] (c2) -- (c3);
\draw[arr] (c3) |- (maint);
\draw[arr] (maint) -| (trim);
\draw[arr] (trim) -- (c1);
\node[below=0.15cm of c2,font=\footnotesize\itshape,text=teal!50!black] {读者可并行 get};
\end{tikzpicture}
\caption{提交三段式与维护反馈：phase2 出锁是读者延迟的关键分界。}
\label{fig:write-e2e}
\end{figure}

\subsection{盘面布局设计图（逻辑）}
\begin{figure}[H]
\centering
\begin{tikzpicture}[
  font=\footnotesize,
  strip/.style={draw=boxline,minimum height=0.7cm,inner sep=2pt,align=center},
  arr/.style={-{Latex},thick}
]
\node[strip,fill=teal!15,minimum width=1.6cm] (sb) {SB\\双槽};
\node[strip,fill=violet!12,minimum width=2.0cm,right=0.15cm of sb] (man) {清单链};
\node[strip,fill=cyan!10,minimum width=3.2cm,right=0.15cm of man] (data) {数据 run 页};
\node[strip,fill=orange!12,minimum width=1.8cm,right=0.15cm of data] (ov) {溢出链};
\node[strip,fill=gray!12,minimum width=1.4cm,right=0.15cm of ov] (free) {free\_};
\node[strip,fill=red!8,minimum width=1.6cm,right=0.15cm of free] (hole) {中部大洞};
\node[above=0.05cm of sb,font=\scriptsize] {低地址};
\node[above=0.05cm of hole,font=\scriptsize] {高水位$\rightarrow$};
\draw[decorate,decoration={brace,amplitude=4pt,mirror},black!50]
  ([yshift=-2pt]data.south west) -- ([yshift=-2pt]ov.south east)
  node[midway,below=6pt,font=\scriptsize] {不可变数据面};
\draw[decorate,decoration={brace,amplitude=4pt},teal!60!black]
  ([yshift=14pt]hole.north west) -- ([yshift=14pt]hole.north east)
  node[midway,above=5pt,font=\scriptsize] {护洞：L0 跳过};
\end{tikzpicture}
\caption{逻辑盘面：控制面（SB/清单）与数据面分离；护洞保护合并释出的中部大洞。}
\label{fig:layout-design}
\end{figure}

\subsection{单元--段--页 设计关系}
\begin{figure}[H]
\centering
\begin{tikzpicture}[
  font=\small,
  c/.style={draw=boxline,rounded corners=2pt,align=left,inner sep=6pt},
  arr/.style={-{Latex},thick,black!50}
]
\node[c,fill=cyan!8] (c1) {\textbf{Cell $[L,R)$}\\Seg$_1$, Seg$_2$, \ldots};
\node[c,fill=orange!8,right=2.2cm of c1] (r1) {\textbf{Run $R_a$}\\seq, ver, fences\\Bloom, sp};
\node[c,fill=violet!8,below=1.0cm of r1] (p1) {\textbf{Pages}\\p0..pn 有序记录};
\draw[arr] (c1) -- node[above,font=\scriptsize]{Seg(键域,页区间)} (r1);
\draw[arr] (r1) -- node[right,font=\scriptsize]{不可变} (p1);
\node[c,fill=gray!6,below=1.0cm of c1] (note) {分裂：收窄键域+页区间\\不搬数据页};
\draw[arr] (c1) -- (note);
\end{tikzpicture}
\caption{单元是视图，run 是物权；分裂零搬移依赖「键域与页区间同时收紧」（坑~1/~20）。}
\label{fig:cell-run}
\end{figure}

\subsection{提交三段式伪代码}
\begin{algorithm}[H]
\caption{Commit（ROUND17 出锁版）}
\begin{algorithmic}[1]
\Require 调用方可未持锁（公开 \texttt{commit}）或已持锁（\texttt{commit\_locked}）
\Ensure 返回时盘上为新提交点，或 \texttt{dirty\_} 待重建
\State $cseq \gets$ \Call{Phase1}{}\Comment{持 XLock：flush\_mem、抓 SB 载荷、发票}
\If{公开路径}
  \State \Call{Unlock}{XLock}
  \State \Call{Phase2Ordered}{$cseq$}\Comment{cmu\_ 按票号：wait\_io；sync；write\_SB；sync}
  \State \Call{Lock}{XLock}
\Else
  \State \Call{Phase2Inline}{}\Comment{维护路径不放锁}
\EndIf
\State \Call{Phase3}{}\Comment{发布 sb\_synced\_、durable\_vers\_；promote\_pending}
\State \Return OK 或 IO（阶段二失败 $\Rightarrow$ dirty\_）
\end{algorithmic}
\end{algorithm}

\begin{algorithm}[H]
\caption{页复用闸（代数）}
\begin{algorithmic}[1]
\Function{CanReusePending}{$gen$}
  \State \Return $gen \le sb\_synced\_$\Comment{\textbf{不是} sb\_built\_：在飞 SB 不算}
\EndFunction
\end{algorithmic}
\end{algorithm}

\subsection{读路径伪代码}
\begin{algorithm}[H]
\caption{Get($k$, $view$)}
\begin{algorithmic}[1]
\State $c \gets$ 目录二分定位含 $k$ 的 cell
\For{Seg $s$ in $c$ 按全序从新到旧}
  \If{Bloom 否定 $k$} \textbf{continue}
  \EndIf
  \State 用 fences/sp 定位页，解码记录
  \If{记录可见于 $view$ 且未被范围墓碑盖住}
    \State \Return 值
  \EndIf
\EndFor
\State \Return NOTFOUND
\end{algorithmic}
\end{algorithm}

\subsection{P2b 分片 + TRIM 伪代码}
\begin{algorithm}[H]
\caption{MergeJob 一片（示意）}
\begin{algorithmic}[1]
\Require 源家族 runs，作业预算页数 $B$
\State $hole \gets$ 最大可用洞（或护洞策略下允许的洞）
\State $part \gets$ \Call{ChoosePart}{$B$, $hole$, frag?}\Comment{可降至 frag\_floor}
\State 归并写出片 $P$（键区间 $[lo,hi)$）至洞或尾
\State 对每个被吃掉前缀的源 $S$：追加 \texttt{M\_RUN\_TRIM}\Comment{锚页 page0 保留}
\State 提交点：清单 + SB；源前缀进入 pending\_free
\end{algorithmic}
\end{algorithm}

\begin{algorithm}[H]
\caption{ChoosePart（含碎片 A）}
\begin{algorithmic}[1]
\State $cap \gets part\_cap(part\_lo)$
\If{$largest\_hole \ge cap$}
  \State \Return 按洞取 $\le cap$ 的片
\ElsIf{$part\_frag$ \textbf{and} $free\ge cap$ \textbf{and} $largest\_hole\ge frag\_floor$}
  \State $part\_frags{+}{+}$
  \State \Return 在 $[frag\_floor, part\_lo)$ 按洞取片
\Else
  \State \Return 落尾（或更大预算）
\EndIf
\end{algorithmic}
\end{algorithm}

\subsection{整理状态机}
\begin{figure}[H]
\centering
\begin{tikzpicture}[
  font=\small,
  st/.style={draw=boxline,rounded corners=3pt,minimum width=2.1cm,minimum height=0.75cm,align=center},
  arr/.style={-{Latex},thick}
]
\node[st,fill=gray!10] (off) {Off};
\node[st,fill=teal!12,right=1.2cm of off] (soft) {Soft};
\node[st,fill=orange!12,right=1.2cm of soft] (hard) {Hard};
\node[st,fill=violet!12,right=1.2cm of hard] (per) {Periodic};
\draw[arr] (off) -- node[above,font=\scriptsize]{configure} (soft);
\draw[arr] (soft) -- node[above,font=\scriptsize]{压力} (hard);
\draw[arr] (hard) -- node[above,font=\scriptsize]{发布前} (per);
\node[below=0.85cm of soft,st,fill=cyan!10,minimum width=3.2cm] (dep) {configure\_deploy\\async+cheap+hole\_fit\\frag+protect+hw};
\node[below=0.85cm of per,st,fill=red!8,minimum width=2.4cm] (pol) {polish\\consolidate+shrink};
\draw[arr] (soft) -- (dep);
\draw[arr] (per) -- (pol);
\end{tikzpicture}
\caption{TidyPolicy：Soft 常驻读友好；Hard/Periodic 步进；polish 重锤。}
\label{fig:tidy-sm}
\end{figure}

\begin{algorithm}[H]
\caption{maybe\_tidy}
\begin{algorithmic}[1]
\If{$policy \in \{Off, Soft\}$} \Return OK
\EndIf
\State $sn \gets$ tidy\_snap()
\If{$sn.runs$ 与 $sn.layout\_pct$ 均未超阈} \Return OK
\EndIf
\If{$policy = Hard$} \Return consolidate\_step()
\Else \Return compact\_file\_step()
\EndIf
\end{algorithmic}
\end{algorithm}

\section{各大踩坑：前后对比详解}
\label{sec:pits}

本章选取对正确性、延迟、峰值影响最大的坑，统一用四段式叙述：
\textbf{现象} $\to$ \textbf{根因} $\to$ \textbf{修前/修后对照}（伪代码或盘面）$\to$ \textbf{验证}。
完整 100+ 条见 \texttt{fcs-pits.html} / DESIGN3 §12。

% ---------- 坑 1 ----------
\subsection{坑 1：幽灵段（切分后 live 翻倍）}
\paragraph{现象}
\texttt{split} 后再合并/recount，同一条记录被计两次；压实产物出现重复副本。

\paragraph{根因}
切分时把父单元段表整份复制给左右子单元；某 run 在一侧键域内\textbf{零记录}仍留下段壳，键域未收窄。

\begin{figure}[H]
\centering
\begin{tikzpicture}[font=\scriptsize,box/.style={draw,rounded corners=1pt,inner sep=3pt,align=center}]
\node[box,fill=oldbg] (o) {修前 R 侧\\Seg(A) 空壳\\键域仍 $[0,100)$};
\node[box,fill=newbg,right=2.5cm of o] (n) {修后 R 侧\\无 Seg(A)\\仅 Seg(B)};
\draw[-{Latex}] (o) -- (n);
\end{tikzpicture}
\caption{幽灵段：空壳把邻居区间记录再次圈入。}
\end{figure}

\begin{fixold}
伪代码（修前）
\begin{lstlisting}[language=pseudo]
function split(C, mid):
  L,R = new cells
  L.segs = C.segs;  R.segs = C.segs   // 整表复制
  // R 上 Seg(A) 可能零记录 → 幽灵
\end{lstlisting}
\end{fixold}

\begin{fixnew}
伪代码（修后）
\begin{lstlisting}[language=pseudo]
function split(C, mid):
  for each seg S in C.segs:
    (klo,khi) = keys_within(S, C.lo, mid)
    if klo < khi: L.add(S narrowed)
    (klo,khi) = keys_within(S, mid, C.hi)
    if klo < khi: R.add(S narrowed)   // 空则丢弃
  retighten(L); retighten(R)          // 坑 20：页区间同步
\end{lstlisting}
\end{fixnew}

\paragraph{验证}
切分 $\to$ 合并 $\to$ verify live $=$ recount；无重复副本。

% ---------- 坑 86 ----------
\subsection{坑 86：写者饥饿（shared\_mutex）}
\paragraph{现象}
\texttt{g7 3 6000 4}：4 读者共享读时写者十几分钟不结束，CPU 近 200\%。

\paragraph{根因}
glibc \texttt{std::shared\_mutex} 读者优先；读者叠加时写者饿死。

\begin{fixold}
\texttt{shared\_lock} 可重入叠加；写者在 \texttt{lock} 上无限等。
\end{fixold}

\begin{fixnew}
自研 \texttt{XLock}：写者优先；读者队列在写者等待时「让一手」；独占持有者重入为空操作。
\end{fixnew}

\paragraph{验证}
g7 多读者 + TSan；写者完成时间回到秒级。

% ---------- 坑 92 ----------
\subsection{坑 92：悬空点墓碑卡住退休}
\paragraph{现象}
全死 run（如 188 页）永不退休；\texttt{exact\_live\_recs} 校准回 live$=1$；verify 报 live 不一致。

\paragraph{根因}
可见性由\textbf{段}定义，但老 \texttt{exact\_live\_recs} 整 run 扫页，把「无更新版本的点墓碑」算成 live。

\begin{figure}[H]
\centering
\begin{tikzpicture}[font=\small,node distance=0.4cm]
\node[draw,fill=oldbg,align=left,inner sep=4pt] (a) {修前：扫全部页\\点墓碑 $\Rightarrow$ live=1\\永不 retire};
\node[draw,fill=newbg,align=left,inner sep=4pt,right=1.8cm of a] (b) {修后：先收集\\仍引用该 run 的 span\\spans$=\emptyset$ $\Rightarrow$ 0};
\end{tikzpicture}
\end{figure}

\begin{fixold}
\begin{lstlisting}[language=pseudo]
exact_live_recs(R):
  n = 0
  for page in R:          // 整条
    for rec in page:
      if not newer_has(rec.key): n++   // 悬空墓碑也算
  return n   // 全死仍可能 = 1
\end{lstlisting}
\end{fixold}

\begin{fixnew}
\begin{lstlisting}[language=pseudo]
exact_live_recs(R):
  spans = segs in cells that reference R
  if spans empty: return 0            // 按定义全死
  count winners only inside spans
\end{lstlisting}
\end{fixnew}

\paragraph{验证}
storm / g9：死 run 退休；verify 干净。

% ---------- 坑 99 / 104 ----------
\subsection{坑 99 / 104：页复用闸与 shrink 提交点}
\paragraph{现象（99）}
NOSPACE / 异步路径下「一页两主」；挂载 CORRUPT。

\paragraph{现象（104）}
\texttt{shrink} 在非提交点截断，掉电丢失\textbf{已提交}数据。

\paragraph{根因}
释放页过早复用（未等 SB 持久）；或维护后 shrink 未先落提交点。ROUND17 起闸比较 \texttt{sb\_synced\_} 而非「已抓载荷」。

\begin{fixold}
pending 释放后，下一次 alloc 即可复用（在飞 SB 未落盘）。\\
\texttt{shrink()} 可在任意维护尾声调用。
\end{fixold}

\begin{fixnew}
\begin{lstlisting}[language=pseudo]
promote_pending:
  only free pages with need_gen <= sb_synced_

shrink:
  require last op was commit point   // maint_commit 先
  then truncate unused high pages
\end{lstlisting}
\end{fixnew}

\paragraph{验证}
g11 t\_nospace；g6 掉电；g14/g15 高水位。

% ---------- 坑 102 ----------
\subsection{坑 102：逻辑版本 u16 回绕}
\paragraph{现象}
第 65536 次逻辑提交后：快照看见「未来」写；钉住 run 暴涨；verify 误报。

\paragraph{根因}
内存 \texttt{Vers=u16} 回绕；盘上本为 32 位。

\begin{fixold}
\texttt{++vers\_} 模 $2^{16}$；快照 view$=65535$ 看见 ver$=0$。
\end{fixold}

\begin{fixnew}
引擎内版本一律 \texttt{u32}（后随 P2 升 u64 盘格式）。
\end{fixnew}

\paragraph{验证}
边界探针 / g12。

% ---------- 坑 112 ----------
\subsection{坑 112：同 seq 家族墓碑复活（ROUND16）}
\paragraph{现象}
键 52000/52002/52004 复活；墓碑与裁过的源同组同 seq。

\paragraph{根因}
老规则假设组内 seq 连续 $\Rightarrow$ $\mathrm{SegKey}<\mathrm{minkey}$ 即更老；同 seq 片打破假设，误丢墓碑。

\begin{figure}[H]
\centering
\begin{tikzpicture}[font=\scriptsize,box/.style={draw,inner sep=3pt,align=left}]
\node[box,fill=oldbg] (o) {T 与 S 同 seq\\older\_may\_contain\\判「下面无副本」\\丢 T $\Rightarrow$ 旧值可见};
\node[box,fill=newbg,right=1.6cm of o] (n) {older\_may\_contain\_grp\\扫 seq$\le$max 且非本组源\\保留 T $\Rightarrow$ 正确删除};
\end{tikzpicture}
\end{figure}

\begin{fixold}
\begin{lstlisting}[language=pseudo]
// 假定组内 seq 连续
if segkey < minkey: return false
\end{lstlisting}
\end{fixold}

\begin{fixnew}
\begin{lstlisting}[language=pseudo]
older_may_contain_grp(k, maxseq):
  for R in runs where R.seq <= maxseq and R not in src_group:
    if key_range_may_contain(R,k) and bloom(R,k): return true
  return false
\end{lstlisting}
\end{fixnew}

\paragraph{验证}
实验头复活 3 键；正式头 0；g16 tomb。

% ---------- r17-1 ----------
\subsection{R17：fsync 出锁（读者 79\,ms $\to$ 0.01\,ms）}
\paragraph{现象}
慢 sync 设备上，提交持锁双 sync 期间 \texttt{get} 平均 $\approx 79\,\mathrm{ms}$。

\paragraph{根因}
可见性本不依赖「已持久」，却把 fsync 放在 XLock 内。

\begin{figure}[H]
\centering
\begin{tikzpicture}[
  font=\scriptsize,
  bar/.style={minimum height=0.45cm,inner sep=1pt,font=\tiny},
  arr/.style={}
]
\node[anchor=west] at (0,1.2) {\textbf{修前}};
\fill[red!40] (1.2,1.1) rectangle (7.2,1.45);
\node at (4.2,1.275) {XLock：flush + sync + SB + sync};
\fill[orange!50] (1.5,0.55) rectangle (7.0,0.9);
\node at (4.25,0.725) {读者 get 等待 $\approx$ 79ms};

\node[anchor=west] at (0,-0.2) {\textbf{修后}};
\fill[orange!30] (1.2,-0.3) rectangle (2.4,0.05);
\node at (1.8,-0.125) {p1};
\fill[violet!35] (2.4,-0.3) rectangle (6.2,0.05);
\node at (4.3,-0.125) {p2 锁外 sync};
\fill[teal!35] (6.2,-0.3) rectangle (7.2,0.05);
\node at (6.7,-0.125) {p3};
\fill[green!40] (2.5,-0.95) rectangle (6.0,-0.6);
\node at (4.25,-0.775) {读者 get $\approx$ 0.01ms};
\end{tikzpicture}
\caption{时间线对照：出锁后读者与阶段二重叠。}
\end{figure}

\paragraph{验证}
\texttt{fsync\_block}；stall；g17；TSan。

% ---------- r17-2 / r18 ----------
\subsection{峰值三刀：hole\_fit / 碎片 A / 并洞 B}
\begin{table}[H]
\centering
\small
\begin{tabular}{@{}llp{5.8cm}@{}}
\toprule
刀 & det200k & 前后要点 \\
\midrule
hole\_fit & $1.96\to\sim1.78\times$ & 中等产物进洞，不再无故落尾 \\
碎片 A & $1.82\to\mathbf{1.71}\times$ & 洞 $<part\_lo$ 时降片下限 \\
护洞 & $\approx$ 不动 & skips$>0$；非 200k KPI \\
并洞 B & $1.71\to1.74$ & 搬挡路抬峰；\textbf{默认关} \\
\bottomrule
\end{tabular}
\caption{养洞相关峰值决策摘要（ROUND17/18 实测）。}
\end{table}

\begin{fixold}
中等产物（$<1024$ 页）一律 append 尾；片下限固定；L0 可吃中部大洞；无并洞。
\end{fixold}

\begin{fixnew}
\begin{lstlisting}[language=pseudo]
alloc_user: if protect and hole>=64 and mid: skip
choose_part:
  if hole_fit or pages>=1024: place in hole
  if frag and hole < part_lo but hole >= floor: lower cap
coalesce: OFF by default   // 1.71→1.74
\end{lstlisting}
\end{fixnew}

% ---------- 家族游标 ----------
\subsection{R17：家族虚拟源游标（CORRUPT 851 $\to$ 0）}
\paragraph{现象}
源片 $>64$ 时 live 账漂移，verify 报数百处 CORRUPT。

\paragraph{根因}
源下标掩码/上限按 64；P2b 片化后同源家族可远超 64。

\begin{fixold}
\texttt{1ull << src\_index} 在 index$\ge$64 未定义；游标数$=$片数。
\end{fixold}

\begin{fixnew}
同 seq 且键域不相交 $\Rightarrow$ 一链一游标（\texttt{open\_chain}/\texttt{next\_run}）；游标数$=$家族数（例 151$\to$8）。
\end{fixnew}

\paragraph{验证}
famw / g17 \texttt{t\_wide\_family}。

% ---------- 小结 ----------
\subsection{对照方法论（可迁移）}
\begin{enumerate}[leftmargin=1.6em]
  \item \textbf{主口径唯一}：读路径、退休、verify 必须共用可见性定义（坑 92/95）。
  \item \textbf{持久点分明}：页复用与 shrink 只能看见已 sync 的 SB（坑 99/104/R17）。
  \item \textbf{假设写进测试}：同 seq 片打破「连续 seq」就要改探测并加判别力用例（坑 112）。
  \item \textbf{KPI 用 A/B 证伪}：护洞叙事正确但 200k 峰不动；并洞机制对但默认关（R18）。
  \item \textbf{锁内只留 CPU}：IO 等待出锁（R17），别用「再加锁」治延迟。
\end{enumerate}

\section{评估与对标}
\label{sec:eval}

\subsection{正确性}
\begin{itemize}[leftmargin=1.6em]
  \item \texttt{make test}：含 f/g/h/s/c/r/b 等数十段；R17/R18 增量 \texttt{g15}--\texttt{g18}。
  \item 掉电：\texttt{g6}/\texttt{h3}（易失缓存 + 随机撕裂）。
  \item 模糊：\texttt{h1} 挂载；\texttt{g5} 交错 API。
  \item 消毒：\texttt{make asan}；关键路径 TSan（w1/g7/g11/g13/g14/g17）。
\end{itemize}
缺口：无小时级 soak；无物理拔电；无多进程文件锁。

\subsection{Peers（FCS / SQLite / LMDB）}
\begin{itemize}[leftmargin=1.6em]
  \item R17 基线旋钮：\texttt{backup/peers-r17/}。
  \item Soft 产品形态（\texttt{configure\_deploy}）：\texttt{backup/peers-soft/}（2026-10-21；\texttt{run\_peers\_soft.sh} / \texttt{run\_all.sh}）。
\end{itemize}
方向性结论（EVAL17 §3.1）：
\begin{itemize}[leftmargin=1.6em]
  \item Soft async+整理：写路径 insert 约 $0.07$--$0.12\times$ SQLite/LMDB；空间 $\sim 2.25\,\mathrm{MB}$ @50k（贴 SQLite、紧于 LMDB $5.54\,\mathrm{MB}$）。
  \item 点查 Soft 下仍强于 SQLite、可优于/追平 LMDB；range 整理后约 $2\times$ LMDB，Soft sync 未整理可达 $\sim 11\times$。
  \item Soft sync 200k/100B 未 polish：insert $1.27\times$ SQLite、range $3.51\times$（与「写换读」一致）。
\end{itemize}
具体 µs/op 以 peers 日志为准；沙箱 2 核抖动常见。

\subsection{峰值 KPI（冻结口径）}
\begin{table}[ht]
\centering
\small
\begin{tabular}{@{}llp{5.5cm}@{}}
\toprule
刀 & det200k 峰 & 决策 \\
\midrule
TRIM + hole\_fit & $\sim 1.75$--$1.81\times$ & 采纳 \\
碎片 A & $\mathbf{\sim 1.71\times}$ & Soft 默认开 \\
护洞 & $\sim 1.81$（相对 1.82） & Soft 开；非 200k KPI \\
并洞 B & $1.71\to 1.74$ & \textbf{默认关} \\
\bottomrule
\end{tabular}
\caption{养洞相关峰值刀（ROUND17/18）。下一刀倾向稀疏产物 / P2c 空间轨。}
\end{table}

\subsection{延迟}
\begin{itemize}[leftmargin=1.6em]
  \item 提交 fsync 出锁：慢盘探针数量级收益；本机快 fsync 上收益有限。
  \item 读者 max 残余常来自维护单元（预算内），而非提交 sync。
  \item 多写者对吞吐加速有限：瓶颈仍在 \texttt{flush\_mem} 等。
\end{itemize}

\section{工程组织与配套材料}

\subsection{目录（产品树）}
\begin{verbatim}
embdb/x/
  include/     fcs.h common.h blockdev.h layout.h store.h(v2)
  tests/       回归 + golden/
  bench/       bench2--7；probe/layout.cpp
  docs/        DESIGN* EVAL17 BATTLE ROUND13--18
  docs/latex/  本文档源
  web/         fcs-anim / fcs-pits 源
  backup/      轮次快照与清扫归档
\end{verbatim}

\subsection{注释模块地图}
\texttt{fcs.h} 文件头含模块 TOC（M01--M07\ldots）；修坏还原见
\texttt{backup/r18/pre-comment-modular/RESTORE.md}。
清扫清单：\texttt{backup/r18/swept/CLEAN\_MANIFEST.md}。

\subsection{可视化}
\begin{itemize}[leftmargin=1.6em]
  \item \texttt{fcs-anim.html}：运行时解剖（React 构建产物）。
  \item \texttt{fcs-pits.html}：坑库；左右示例数据动画 + 深析；含 R15--R18 扩展批。
  \item 重建：\texttt{python3 web/pits/build.py}；\texttt{cd web \&\& npm run build}。
\end{itemize}

\subsection{已知工程债}
\begin{itemize}[leftmargin=1.6em]
  \item 单头文件体量（约 $8\,\mathrm{k}$ 行级）：拆分与本地 git 流程仍弱。
  \item 旋钮多：部署请走 \texttt{configure\_deploy}/\texttt{TidyPolicy}，避免漏配。
  \item P2c（\texttt{M\_RUN\_HOLE}）仅占 op 号；再压峰走空间轨而非更多 cliff 旋钮。
\end{itemize}


\appendix
\section{公开 API 速查（部署相关）}
\label{app:api}

下列为部署与整理相关入口摘要；完整签名以 \texttt{include/fcs.h} 为准。

\begin{lstlisting}[language=fcsapi,caption={整理与养洞旋钮}]
void configure_deploy(bool async = true, u32 hw_pct = 150);
RC   polish();
void set_tidy_policy(TidyPolicy p);
TidySnap tidy_snap() const;
RC   maybe_tidy();

void set_part_hole_fit(bool on);
void set_part_frag(bool on);
void set_part_frag_floor(u32 n);
void set_part_coalesce(bool on);      // 默认建议关
void set_protect_merge_holes(bool on);

void set_async_writes(bool on);
void set_merge_cheap_pages(u32 n);
void set_hw_reclaim_pct(u32 pct);
void set_maint_mode(int mode);        // 0 / 1 / 3
u64  durable_vers() const;
\end{lstlisting}

\begin{description}[leftmargin=1.4em]
  \item[\texttt{configure\_deploy}] Soft 一揽子：async、cheap、hole\_fit、frag~A、护洞、hw；必要时把 maint 从 0 提到 1。
  \item[\texttt{polish}] 全库归一 + 尾部回落；写放大高，用于发布前。
  \item[\texttt{maybe\_tidy}] Soft 直接 OK；Hard/Periodic 在 runs 或 layout\% 超阈时步进。
\end{description}

\section{复现命令}
\label{app:repro}

均在 \texttt{embdb/x} 下、WSL 环境：

\begin{lstlisting}[language=bash,caption={回归与关键探针}]
make -j2 -s build/g18 && ./build/g18
make -j2 -s test          # 全量较久
# layout 峰值（需先编 probe）
g++ -std=c++17 -O2 -pthread -Iinclude \
  bench/probe/layout.cpp -o build/layout
# det 口径示例：budget=0 units=8
./build/layout 200000 500 32 1 1 0 0 0 0 8
# 护洞 / 碎片：见 backup/r18/run_layout_*.sh
bash backup/peers-r17/run_peers_r17.sh      # R17 基线
bash backup/peers-soft/run_peers_soft.sh    # Soft 产品形态
# 或：bash backup/peers-soft/run_all.sh     # make test + Soft peers
\end{lstlisting}

\begin{lstlisting}[language=bash,caption={本文档编译}]
cd docs/latex
make          # xelatex ×2
# 产物：build/main.pdf
\end{lstlisting}

\subsection*{参考文献（仓库内）}
\begin{itemize}[leftmargin=1.6em]
  \item \texttt{docs/README.md} — Markdown 分类索引
  \item \texttt{docs/10-design/DESIGN3.md} — 设计与历史坑清单 §12
  \item \texttt{docs/00-entry/EVAL17.md} — 冻结全量评估
  \item \texttt{docs/20-rounds/ROUND15.md}--\texttt{docs/20-rounds/ROUND18.md} — 各轮执行与数字
  \item \texttt{docs/30-discuss/DISCUSS15\_FSYNC.md} — fsync 出锁讨论稿
  \item \texttt{docs/10-design/P2\_DESIGN.md} — 格式 v2 设计
  \item \texttt{docs/00-entry/BATTLE.md} — 升格与分叉对照
  \item \texttt{backup/README.md} — 归档 / peers 日志分类
\end{itemize}


\end{document}
